% ch05_structures.tex — Structures de Donn\'ees
\chapter{Structures de Donn\'ees}
\label{ch:structures}

\begin{intuition}
Un scientifique ne travaille jamais avec une seule donn\'ee~: il manipule des s\'eries de
mesures, des tableaux de r\'esultats, des catalogues d'esp\`eces. Python offre plusieurs
\emph{structures de donn\'ees} pour organiser ces informations --- chacune avec ses forces,
comme les diff\'erents contenants d'un laboratoire (\'eprouvettes, bo\^ites de P\'etri,
classeurs).
\end{intuition}

% ============================================================
\section{Les listes}
\index{liste}

\begin{definition}[Liste]
Une \textbf{liste} est une collection \emph{ordonn\'ee} et \emph{modifiable} d'\'el\'ements.
On la cr\'ee avec des crochets~: \py{[e1, e2, ...]}.
\end{definition}

\subsection{Cr\'eation et indexation}
\index{indexation}

\begin{codeblock}[title=\textbf{Python}]
\begin{minted}{python}
# Mesures de température (°C) sur une semaine
temperatures = [18.2, 19.5, 17.8, 21.0, 22.3, 20.1, 19.7]

print(f"Premier jour  : {temperatures[0]}°C")
print(f"Dernier jour  : {temperatures[-1]}°C")
print(f"Nombre de mesures : {len(temperatures)}")
\end{minted}
\end{codeblock}

\begin{codeoutput}
\begin{verbatim}
Premier jour  : 18.2°C
Dernier jour  : 19.7°C
Nombre de mesures : 7
\end{verbatim}
\end{codeoutput}

\subsection{Slicing (d\'ecoupage)}
\index{slicing}

\begin{codeblock}[title=\textbf{Python}]
\begin{minted}{python}
# Jours ouvrés (lundi à vendredi)
jours_ouvres = temperatures[0:5]
print(f"Jours ouvrés : {jours_ouvres}")

# Un élément sur deux
print(f"Un sur deux : {temperatures[::2]}")

# Ordre inverse
print(f"Inversé : {temperatures[::-1]}")
\end{minted}
\end{codeblock}

\begin{codeoutput}
\begin{verbatim}
Jours ouvrés : [18.2, 19.5, 17.8, 21.0, 22.3]
Un sur deux : [18.2, 17.8, 22.3, 19.7]
Inversé : [19.7, 20.1, 22.3, 21.0, 17.8, 19.5, 18.2]
\end{verbatim}
\end{codeoutput}

\subsection{M\'ethodes utiles}
\index{append}\index{pop}\index{sort}

\begin{codeblock}[title=\textbf{Python}]
\begin{minted}{python}
masses = [12.5, 8.3, 15.7, 10.2]

masses.append(9.8)       # ajouter à la fin
print(f"Après append : {masses}")

dernier = masses.pop()   # retirer le dernier
print(f"Retiré : {dernier}, liste : {masses}")

masses.sort()            # trier en place
print(f"Trié : {masses}")

masses.reverse()         # inverser en place
print(f"Inversé : {masses}")
\end{minted}
\end{codeblock}

\begin{codeoutput}
\begin{verbatim}
Après append : [12.5, 8.3, 15.7, 10.2, 9.8]
Retiré : 9.8, liste : [12.5, 8.3, 15.7, 10.2]
Trié : [8.3, 10.2, 12.5, 15.7]
Inversé : [15.7, 12.5, 10.2, 8.3]
\end{verbatim}
\end{codeoutput}

% ============================================================
\section{Compr\'ehensions de listes}
\index{compr\'ehension de liste}

Les compr\'ehensions offrent une syntaxe compacte pour cr\'eer des listes.

\begin{codeblock}[title=\textbf{Python}]
\begin{minted}{python}
# Carrés des entiers de 1 à 10
carres = [x**2 for x in range(1, 11)]
print(f"Carrés : {carres}")

# Filtrer les températures > 20°C
temperatures = [18.2, 19.5, 17.8, 21.0, 22.3, 20.1, 19.7]
chaudes = [t for t in temperatures if t > 20]
print(f"Jours chauds : {chaudes}")

# Convertir en Fahrenheit
fahrenheit = [t * 9/5 + 32 for t in temperatures]
print(f"Fahrenheit : {[f'{f:.1f}' for f in fahrenheit]}")
\end{minted}
\end{codeblock}

\begin{codeoutput}
\begin{verbatim}
Carrés : [1, 4, 9, 16, 25, 36, 49, 64, 81, 100]
Jours chauds : [21.0, 22.3, 20.1]
Fahrenheit : ['64.8', '67.1', '64.0', '69.8', '72.1', '68.2', '67.5']
\end{verbatim}
\end{codeoutput}

% ============================================================
\section{Les tuples}
\index{tuple}

\begin{definition}[Tuple]
Un \textbf{tuple} est une collection \emph{ordonn\'ee} mais \emph{immuable}
(non modifiable). On le cr\'ee avec des parenth\`eses~: \py{(e1, e2, ...)}.
\end{definition}

\begin{codeblock}[title=\textbf{Python}]
\begin{minted}{python}
# Constantes physiques (ne doivent pas changer)
constantes = ("vitesse lumière", 299_792_458, "m/s")
print(f"{constantes[0]} = {constantes[1]} {constantes[2]}")

# Déballage (unpacking)
nom, valeur, unite = constantes
print(f"{nom} : {valeur} {unite}")

# Tentative de modification → erreur
# constantes[1] = 300_000_000  # TypeError !
\end{minted}
\end{codeblock}

\begin{codeoutput}
\begin{verbatim}
vitesse lumière = 299792458 m/s
vitesse lumière : 299792458 m/s
\end{verbatim}
\end{codeoutput}

\begin{remarque}
Utilisez les tuples pour des donn\'ees qui ne doivent pas \^etre modifi\'ees~:
coordonn\'ees $(x, y)$, constantes physiques, cl\'es de dictionnaire.
\end{remarque}

% ============================================================
\section{Les dictionnaires}
\index{dictionnaire}

\begin{definition}[Dictionnaire]
Un \textbf{dictionnaire} associe des \emph{cl\'es} \`a des \emph{valeurs}.
On le cr\'ee avec des accolades~: \py{\{cle: valeur, ...\}}.
\end{definition}

\subsection{Le tableau p\'eriodique comme dictionnaire}
\index{tableau p\'eriodique}

\begin{codeblock}[title=\textbf{Python}]
\begin{minted}{python}
elements = {
    "H":  {"nom": "Hydrogène",  "Z": 1,  "masse": 1.008},
    "He": {"nom": "Hélium",     "Z": 2,  "masse": 4.003},
    "C":  {"nom": "Carbone",    "Z": 6,  "masse": 12.011},
    "N":  {"nom": "Azote",      "Z": 7,  "masse": 14.007},
    "O":  {"nom": "Oxygène",    "Z": 8,  "masse": 15.999},
    "Fe": {"nom": "Fer",        "Z": 26, "masse": 55.845},
}

# Accès à un élément
carbone = elements["C"]
print(f"Carbone : Z={carbone['Z']}, masse={carbone['masse']} u")

# Masse molaire de l'eau H2O
masse_eau = 2 * elements["H"]["masse"] + elements["O"]["masse"]
print(f"Masse molaire H₂O = {masse_eau:.3f} g/mol")
\end{minted}
\end{codeblock}

\begin{codeoutput}
\begin{verbatim}
Carbone : Z=6, masse=12.011 u
Masse molaire H₂O = 18.015 g/mol
\end{verbatim}
\end{codeoutput}

\subsection{It\'eration sur un dictionnaire}

\begin{codeblock}[title=\textbf{Python}]
\begin{minted}{python}
# Afficher tous les éléments
for symbole, info in elements.items():
    print(f"{symbole:>2} : {info['nom']:<12} Z={info['Z']:>2}")
\end{minted}
\end{codeblock}

\begin{codeoutput}
\begin{verbatim}
 H : Hydrogène    Z= 1
He : Hélium       Z= 2
 C : Carbone      Z= 6
 N : Azote        Z= 7
 O : Oxygène      Z= 8
Fe : Fer          Z=26
\end{verbatim}
\end{codeoutput}

% ============================================================
\section{Les ensembles}
\index{ensemble}\index{set}

\begin{codeblock}[title=\textbf{Python}]
\begin{minted}{python}
# Éléments détectés dans deux échantillons
echantillon_A = {"H", "C", "O", "N", "S"}
echantillon_B = {"H", "C", "O", "Fe", "Ca"}

print(f"Communs    : {echantillon_A & echantillon_B}")
print(f"Tous       : {echantillon_A | echantillon_B}")
print(f"Unique à A : {echantillon_A - echantillon_B}")
\end{minted}
\end{codeblock}

\begin{codeoutput}
\begin{verbatim}
Communs    : {'O', 'H', 'C'}
Tous       : {'S', 'O', 'Fe', 'H', 'Ca', 'N', 'C'}
Unique à A : {'S', 'N'}
\end{verbatim}
\end{codeoutput}

% ============================================================
% TikZ comparison
\begin{figure}[ht]
\centering
\begin{tikzpicture}[
    header/.style={rectangle, draw, fill=blue!60, text=white,
        minimum width=3.5cm, minimum height=0.7cm, font=\bfseries, align=center},
    cell/.style={rectangle, draw, minimum width=3.5cm, minimum height=0.6cm,
        font=\small, align=center},
]
% Liste
\node[header] (lh) {Liste};
\node[cell, below=0pt of lh] (l1) {Ordonn\'ee};
\node[cell, below=0pt of l1] (l2) {Modifiable};
\node[cell, below=0pt of l2] (l3) {\py{[1, 2, 3]}};

% Tuple
\node[header, right=0.5cm of lh] (th) {Tuple};
\node[cell, below=0pt of th] (t1) {Ordonn\'e};
\node[cell, below=0pt of t1] (t2) {Immuable};
\node[cell, below=0pt of t2] (t3) {\py{(1, 2, 3)}};

% Dict
\node[header, right=0.5cm of th] (dh) {Dictionnaire};
\node[cell, below=0pt of dh] (d1) {Cl\'e $\to$ Valeur};
\node[cell, below=0pt of d1] (d2) {Modifiable};
\node[cell, below=0pt of d2] (d3) {\py{\{"a": 1\}}};

% Set
\node[header, right=0.5cm of dh] (sh) {Ensemble};
\node[cell, below=0pt of sh] (s1) {Non ordonn\'e};
\node[cell, below=0pt of s1] (s2) {Pas de doublons};
\node[cell, below=0pt of s2] (s3) {\py{\{1, 2, 3\}}};
\end{tikzpicture}
\caption{Comparaison des structures de donn\'ees en Python.}
\label{fig:structures-comparaison}
\end{figure}

% ============================================================
\section{Erreurs courantes}

\begin{errmark}
\begin{enumerate}
    \item \textbf{Index hors limites (\py{IndexError})}~: une liste de 5 \'el\'ements a
    des indices de 0 \`a 4. Acc\'eder \`a l'indice 5 provoque une erreur.
    \begin{minted}{python}
mesures = [1.2, 3.4, 5.6]
print(mesures[3])  # IndexError !
# Correct : mesures[2] ou mesures[-1]
    \end{minted}

    \item \textbf{Modifier une liste pendant l'it\'eration}~: supprimer des \'el\'ements
    d'une liste en la parcourant produit des r\'esultats impr\'evisibles.
    \begin{minted}{python}
# ERREUR : saute des éléments
donnees = [1, -2, 3, -4, 5]
for d in donnees:
    if d < 0:
        donnees.remove(d)  # Dangereux !

# CORRECT : créer une nouvelle liste
donnees = [d for d in donnees if d >= 0]
    \end{minted}

    \item \textbf{Cl\'e inexistante (\py{KeyError})}~: acc\'eder \`a une cl\'e absente d'un
    dictionnaire. Utilisez \py{dict.get(cle, defaut)} pour \'eviter l'erreur.

    \item \textbf{Confondre liste et tuple}~: tenter de modifier un tuple avec
    \py{t[0] = 5} provoque un \py{TypeError}.
\end{enumerate}
\end{errmark}

% ============================================================
\section{Mini-projet~: Gestion des notes \'etudiantes}

\begin{projet}
Cr\'eez un syst\`eme de gestion de notes utilisant un dictionnaire~:
\begin{itemize}
    \item Cl\'es~: noms des \'etudiants
    \item Valeurs~: listes de notes
    \item Fonctions~: ajouter un \'etudiant, ajouter une note, calculer la moyenne,
          trouver le major de promotion
\end{itemize}
\end{projet}

\begin{codeblock}[title=\textbf{Python}]
\begin{minted}{python}
def creer_promotion():
    """Crée un dictionnaire vide de promotion."""
    return {}

def ajouter_etudiant(promo, nom):
    """Ajoute un étudiant avec une liste de notes vide."""
    promo[nom] = []

def ajouter_note(promo, nom, note):
    """Ajoute une note à un étudiant."""
    promo[nom].append(note)

def moyenne(notes):
    """Calcule la moyenne d'une liste de notes."""
    return sum(notes) / len(notes) if notes else 0

def afficher_resultats(promo):
    """Affiche les résultats de la promotion."""
    print(f"{'Étudiant':<15} {'Notes':<25} {'Moyenne':>8}")
    print("-" * 50)
    for nom, notes in promo.items():
        notes_str = ", ".join(f"{n:.1f}" for n in notes)
        moy = moyenne(notes)
        print(f"{nom:<15} {notes_str:<25} {moy:>7.2f}")

# Utilisation
promo = creer_promotion()
for nom in ["Alice", "Bob", "Claire"]:
    ajouter_etudiant(promo, nom)

ajouter_note(promo, "Alice", 15.5)
ajouter_note(promo, "Alice", 17.0)
ajouter_note(promo, "Alice", 14.0)
ajouter_note(promo, "Bob", 12.0)
ajouter_note(promo, "Bob", 11.5)
ajouter_note(promo, "Bob", 13.0)
ajouter_note(promo, "Claire", 18.0)
ajouter_note(promo, "Claire", 16.5)
ajouter_note(promo, "Claire", 19.0)

afficher_resultats(promo)

major = max(promo, key=lambda nom: moyenne(promo[nom]))
print(f"\nMajor de promotion : {major} ({moyenne(promo[major]):.2f}/20)")
\end{minted}
\end{codeblock}

\begin{codeoutput}
\begin{verbatim}
Étudiant        Notes                     Moyenne
--------------------------------------------------
Alice           15.5, 17.0, 14.0            15.50
Bob             12.0, 11.5, 13.0            12.17
Claire          18.0, 16.5, 19.0            17.83

Major de promotion : Claire (17.83/20)
\end{verbatim}
\end{codeoutput}

% ============================================================
\section{Exercices}

\begin{exercice}[$\star$ --- Manipulation de listes]
Cr\'eez une liste de 10 mesures de pH. Calculez la moyenne, le minimum et le maximum
\emph{sans utiliser} les fonctions \py{sum}, \py{min}, \py{max}.
\end{exercice}

\begin{exercice}[$\star$ --- Filtrage]
\`A partir d'une liste de temp\'eratures, cr\'eez avec une compr\'ehension de liste
une nouvelle liste ne contenant que les valeurs entre 15\,\textdegree C et 25\,\textdegree C.
\end{exercice}

\begin{exercice}[$\star\star$ --- Compteur de fr\'equences]
\'Ecrivez une fonction qui re\c{c}oit une cha\^ine de caract\`eres et renvoie un dictionnaire
comptant la fr\'equence de chaque lettre. Testez avec une s\'equence ADN
(<<~ATCGGATCCA~>>).
\end{exercice}

\begin{exercice}[$\star\star$ --- Matrice comme liste de listes]
Repr\'esentez une matrice $3\times 3$ comme une liste de listes. \'Ecrivez une fonction
qui calcule la transpos\'ee. V\'erifiez avec la matrice identit\'e.
\end{exercice}

\begin{exercice}[$\star\star\star$ --- Analyse de s\'equence prot\'eique]
Cr\'eez un dictionnaire associant chaque acide amin\'e \`a sa masse molaire.
\'Ecrivez une fonction qui, \'etant donn\'ee une s\'equence d'acides amin\'es (lettres),
calcule la masse totale de la prot\'eine. Testez avec <<~MKVLWA~>>.
\end{exercice}

% ============================================================
\begin{formules}{R\'esum\'e du chapitre~: Structures de donn\'ees}
\begin{itemize}
    \item \textbf{Liste} \py{[...]}~: ordonn\'ee, modifiable, m\'ethodes \py{append},
    \py{pop}, \py{sort}
    \item \textbf{Compr\'ehension}~: \py{[expr for x in iterable if cond]}
    \item \textbf{Tuple} \py{(...)}~: ordonn\'e, immuable, id\'eal pour les constantes
    \item \textbf{Dictionnaire} \py{\{cle: val\}}~: acc\`es par cl\'e, m\'ethodes
    \py{.keys()}, \py{.values()}, \py{.items()}
    \item \textbf{Ensemble} \py{\{...\}}~: pas de doublons, op\'erations $\cap$, $\cup$, $\setminus$
    \item \textbf{Indexation}~: commence \`a 0, indices n\'egatifs depuis la fin
    \item \textbf{Slicing}~: \py{liste[debut:fin:pas]}
\end{itemize}
\end{formules}
